\begin{question}
\noindent A historian is studying an ancient astronomical diagram carved into a stone tablet, used to track the position of stars. The diagram shows a straight line $AB$ with a point $C$ on it, and a ray $CD$ representing a sight-line to a star.

\begin{center}
\begin{tikzpicture}[scale=1]
    % Define the points
    \coordinate (A) at (0,0);
    \coordinate (B) at (4,0);
    \coordinate (C) at (2,0);
    \coordinate (D) at ($(C)+(138:2cm)$);

    % Draw the straight line and the ray
    \draw (A) -- (B);
    \draw (C) -- (D);

    % Draw the angles
    \pic [draw, "$138^{\circ}$", angle eccentricity=1.5, angle radius=0.6cm] {angle = B--C--D};
    \pic [draw, "$x^{\circ}$", angle eccentricity=1.8, angle radius=0.8cm] {angle = D--C--A};

    % Place the labels
    \node at (A) [left,font=\small] {A};
    \node at (B) [right,font=\small] {B};
    \node at (C) [below,font=\small] {C};
    \node at (D) [above,font=\small] {D};
\end{tikzpicture}
\end{center}

\noindent $ACB$ is a straight line.

\noindent Work out the value of $x$.

\vspace{4cm}

\answereq{x}{2}

\end{question}